
\chapter{Trabajo final realizado}
\section{Puntos Función finales del proyecto}

Con respecto a la entrega intermedia se han mantenido en el mismo estado los ILF y los EIF.

\begin{figure}[H]
  \begin{center}
    \includegraphics[width=1\textwidth]{PF-FINAL/ILF.png}
  \end{center}
  \caption{Puntos función de los ILF y el estado de los mismos.}
  \label{fig:ILF-final}
\end{figure}


\begin{figure}[H]
  \begin{center}
    \includegraphics[width=1\textwidth]{PF-FINAL/EIF.png}
  \end{center}
  \caption{Puntos función de los EIF y el estado de los mismos.}
  \label{fig:EIF-final}
\end{figure}


\begin{figure}[H]
  \begin{center}
    \includegraphics[width=1\textwidth]{PF-FINAL/EI.png}
  \end{center}
  \caption{Puntos función de los EI y el estado de los mismos.}
  \label{fig:EI-final}
\end{figure}

\begin{figure}[H]
  \begin{center}
    \includegraphics[width=1\textwidth]{PF-FINAL/EO.png}
  \end{center}
  \caption{Puntos función de los EO y el estado de los mismos.}
  \label{fig:EO-final}
\end{figure}

\begin{figure}[H]
  \begin{center}
    \includegraphics[width=1\textwidth]{PF-FINAL/EQ.png}
  \end{center}
  \caption{Puntos función de los EQ y el estado de los mismos.}
  \label{fig:EQ-final}
\end{figure}


\section{Resumen del trabajo realizado para la entrega final}
% En la última parte del proyecto, se han finalizado los módulos restantes para la finalización de la formación en el programa del TFG. En el ámbito del de diseño y desarrollo se han implementado las pantallas de calendario, que permitirán al personal investigador la reserva de horas de cómputo. En el requisito de de la prevención de acceso nativo junto a la herramienta de code-server se ha diseñado la solución final que junto scripts son capaces de proporcionar entornos de desarrollo virtualizados y seguros. Además de esto, se ha perfeccionado la integración de terceros de gitlab, en el ámbito de armonización de los datos en nuestra base de datos de verdad. Por último se han creado pruebas de usabilidad del producto mínimo viable desarrollado así como su evaluación mediante el UEQ. Se ha realizado el análisis de los resultados.

% En la última etapa del proyecto se han completado los módulos restantes que conforman el alcance del TFG. En el ámbito del diseño y desarrollo, se han implementado las pantallas de calendario, que permitirán al personal investigador reservar franjas horarias de cómputo de forma autónoma. Respecto a la prevención del acceso nativo a las máquinas, se ha diseñado e implementado la solución definitiva que, mediante scripts especializados, es capaz de proporcionar entornos de desarrollo virtualizados y seguros en combinación con la herramienta code-server. Asimismo, se ha perfeccionado la integración con GitLab, consolidando la armonización y consistencia de los datos en la base de datos como fuente de verdad del sistema. Por último, se han diseñado y llevado a cabo pruebas de usabilidad sobre el producto mínimo viable desarrollado, evaluando los resultados obtenidos mediante el cuestionario UEQ (User Experience Questionnaire) y analizando las conclusiones derivadas de dicha evaluación.

En la última etapa del proyecto se han completado los módulos restantes de formación del TFG. En el ámbito de diseño de la solución y su desarrollo se han implementado las pantallas de calendario, que permitirán al personal investigador la reserva de franjas horarios de cómputo. Esto, permite una mayor organización de los recursos de las máquinas. Respecto al acceso nativo de las máquinas, se ha diseñado e implementado una solución que mediante unos scripts de automatización es capaz de proporcionar entornos de desarrollo virtualizados y seguros junto a los contenedores de code-server, se ha pasado a producción en uno de los servidores del laboratorio. Además, se ha mejorado la integración con el servicio de Gitlab, consolidando la armonización y consistencia de los datos presentes en la base de datos como fuente única de verdad en el sistema. Por último se han diseñado y llevado a cabo pruebas de usabilidad sobre el producto mínimo viable para su posterior evaluación mediante técnicas de evaluación en entrevistas y el relleno del cuestionario del UEQ para su posterior análisis. Con respecto al resto de integraciones, se realiza un pequeño marco teórico para poder conocer que recursos están a disposición para el desarrollo e incorporación. Debido a la mejora e implementación de la integración de GitLab no se ha podido llevar a cabo la integración de NUKI. 


En la tabla \ref{fig:RESUMEN-FINAL} se observa la cantidad de horas destinadas, cuanto tiempo por Punto de Función y qué porcentaje del proyecto se ha realizado. En la tabla \ref{fig:RESUMEN-FINAL} se puede observar la cantidad de horas destinadas para cada tarea y en el diagrama \ref{fig:graficoAnalisisPlan-final} una representación visual. 

\begin{figure}[H]
  \begin{center}
    
\includegraphics[width=1\textwidth]{PF-FINAL/RESUMEN.png}
  \end{center}
  \caption{Tabla Resumen de Estado del Proyecto.}
  \label{fig:RESUMEN-FINAL}
\end{figure}


\begin{figure}[H]
  \begin{center}
    \includegraphics[width=1\textwidth]{PF-FINAL/horas_tareas.png}
  \end{center}
  \caption{Tabla de Análisis de horas destinadas al proyecto.}
  \label{fig:AnalisisPlan-final}
\end{figure}


\begin{figure}[H]
  \begin{center}
    \includegraphics[width=1\textwidth]{PF-FINAL/horas_tareas_grafico.png}
  \end{center}
  \caption{Gráfico visual de las horas destinadas.}
  \label{fig:graficoAnalisisPlan-final}
\end{figure}


\section{Cronograma final de Gannt}
En la figura \ref{fig:gantt_mermaid_final} se puede ver el cronograma final de desarrollo.

\clearpage
\begin{sidewaysfigure}[pH]
    \centering
    \includegraphics[width=\textheight, keepaspectratio]{gannt-final/gannt_final.png}
    \caption{Cronograma del desarrollo final.}
    \label{fig:gantt_mermaid_final}
\end{sidewaysfigure}


